% ============================================================================
%  Declaración de uso de IA del Subdocumento 13 (Comunicado 10, sección 7.2).
%  \usoIA{sección}{herramienta}{finalidad}{nivel en texto}{nivel en diagramas}{revisión humana}
%  La revisión humana la firma quien revisó, con nombre y cargo en audIT.
% ============================================================================
\usoIA{Introducción}{Claude Opus 5.5 en Claude Code}{Redacción inicial de la estrategia y de la cartera, y diagrama de la cartera generado con un guion desde el modelo definido por el equipo}{Alto}{Medio}{Ignacio Vera, Arquitecto de Solución: problema de cada innovación, lugar en la arquitectura y figura de la cartera.}
\usoIA{13.1, 13.4 y 13.5}{Claude Opus 5.5 en Claude Code}{Ordenamiento del texto por los siete elementos del Artículo 29 y verificación de las citas al Caso}{Medio}{Ninguno}{Ignacio Vera, Arquitecto de Solución: siete elementos del Artículo 29 y citas al Caso.}
\usoIA{13.2 y ficha 2 del Formulario T-19}{OpenAI Codex (GPT-6)}{Diseño de la programación de intervenciones por oportunidad de patio, indicadores, riesgos y gráfico}{Alto}{Alto}{Naomi Núñez, Líder de Calidad, e Ignacio Vera y Marcel Gutiérrez, Arquitectos de Solución: proceso de programación por oportunidad de patio, indicadores, riesgos y figura.}
\usoIA{13.2 y ficha 2 del Formulario T-19}{OpenAI Codex}{Cierre de la flota propia en el mes 9 y calendario del Subdocumento 7}{Alto}{Ninguno}{Naomi Núñez, Líder de Calidad, e Ignacio Vera y Marcel Gutiérrez, Arquitectos de Solución: calendario de la innovación 2 frente al Subdocumento 7.}
\usoIA{13.3 y ficha 3 del Formulario T-19}{Claude Opus 5.5 en Claude Code}{Búsqueda de fichas de fabricante, redacción inicial desde las decisiones del equipo, y actualización para la oferta final: el kilometraje de los semirremolques pasa al criterio 26 y el beneficio queda en dos partes medidas por separado, sin valor monetario en el escenario base}{Alto}{Medio}{Naomi Núñez, Líder de Calidad, e Ignacio Vera y Marcel Gutiérrez, Arquitectos de Solución: delimitación frente al criterio 26, indicadores, inversión y costo coherentes con el Sobre N.º 3, y ficha 3 del Formulario T-19.}
\usoIA{Formulario T-19}{Claude Opus 5.5 en Claude Code}{Armado de las cinco fichas desde las secciones 13.1 a 13.5}{Alto}{Ninguno}{Naomi Núñez, Líder de Calidad, e Ignacio Vera y Marcel Gutiérrez, Arquitectos de Solución: correspondencia de las cinco fichas con las secciones 13.1 a 13.5.}

\usoIA{13.5 y ficha T-19 tipo 5}{Asistente LLM y OpenAI Codex}{Integración del texto de experiencia de usuario; depuración de madurez y calendario; definición del indicador de recomendación y alcance del beneficio}{Alto}{Ninguno}{Naomi Núñez, Líder de Calidad, e Ignacio Vera y Marcel Gutiérrez, Arquitectos de Solución: claridad y consistencia de la innovación 5, madurez actual y objetivos, ventanas de validación y producción, fórmula de efectividad y distinción entre metas y resultados.}
\usoIA{13.2 y ficha 2}{Codex}{Conciliación de lotes, piloto y medición de beneficio}{Alto}{Ninguno}{Naomi Núñez, Líder de Calidad, e Ignacio Vera y Marcel Gutiérrez, Arquitectos de Solución: lotes, piloto y medición del beneficio frente al Subdocumento 7 y a la ficha 2.}
\usoIA{13.1 a 13.5 y Formulario T-19}{Claude Opus 5.5 en Claude Code}{Conciliación de los riesgos de adopción con el Formulario T-16 y exposición residual por ficha, delimitación de la innovación 5 frente al criterio 28, nombres de las innovaciones 3 y 5 y meses de montaje alineados con el Formulario T-14, impacto económico de la cartera según el presupuesto del Sobre N.º 3 y diagrama de la cartera regenerado con su guion}{Alto}{Medio}{Naomi Núñez, Líder de Calidad, e Ignacio Vera y Marcel Gutiérrez, Arquitectos de Solución: exposición inicial y residual de cada riesgo frente al Formulario T-16, alcance base y mérito innovador de las fichas 3 y 5, nombres y meses frente al Formulario T-14 y ausencia de horas y montos.}
